\chapter{前人的工作基础}\label{chap:introduction}

本章首先详细介绍整个系统所使用的两个前人的研究成果，分别是流式模糊哈希和基于流式模糊哈希的相似性查找系统。接着介绍使用这些研究成果的原因和各自的优点。

\section{使用流式模糊哈希的原因及其优点}
在网络流量中识别重复的音视频文件，需要一个基准来识别文件是否相同。根据\cite{aoli2010}的论文，一般使用MD5对文件是否相同进行检测。MD5是一种很常见的哈希算法，哈希算法从数学上看是一种函数映射，其可以将任意大小的文件或者任意长度的字符串映射成一个固定长度的字符串。

但在真实的网络环境中，由于人们的观看习惯，例如拖动在线视频的进度条，或者手动关闭视频网页，导致网络中传输的音视频文件常常会出现缺损，因此，对于同一个音视频文件版本，很难出现两个完全相同的文件，往往需要对文件的相似性进行检测。

普通哈希算法具有雪崩效应。当文件中的任意一个比特位发生改变时，会导致生成的哈希值完全不同。虽然这一特性对于检测文件是否缺损，或者检测文件是否改变有着重要作用，但却无法对于两个文件进行相似性检测。

为了能够对文件的相似性进行检测，提出了模糊哈希算法。和普通哈希相比较，模糊哈希可以比较两个不同版本文件的相似度，其主要使用在计算机数字取证和恶意代码分析领域。模糊哈希的特点如下：

非传播（Non-propagation）：只有与二进制的变化部分线性对应的模糊哈希值部分才会被更改。这意味着文件中任何部分的微小变化都不会引起哈希值急剧变化。

对齐鲁棒性（Alignment robustness）：大多数哈希算法都对两个文件是否对齐十分敏感。仅仅向文件中删除或插入单个字节将导致最后生成两个完全不同的哈希值。模糊哈希算法的核心是用于在二进制中生成一系列重置点的弱哈希（rolling hash）。且该重置点仅取决于直接上下文。

下面简单介绍一下生成模糊哈希的步骤，如公式~\ref{eq:chufa}~所示，模糊哈希使用名为块大小b的变量来触发重置点，块大小可以用1计算，这可以确保某个文件的模糊哈希结果长度既不会太长不能比较，也不会太短从而避免冲突。 bmin是常量最小块大小，S是常量预期模糊散列长度，n是输入文件大小。

\begin{equation} \label{eq:chufa}
b_{init}=b_{min}2^{\lfloor \log_2(\frac{n}{Sb_{min}}) \rfloor}
\end{equation}

k窗口的弱哈希函数在k字节的输入序列：$c_1$,$c_2$,... ,$c_k$，当满足公式~\ref{eq:rollinghash}~时，重置点将位于$c_k$。根据统计上，当b值越小时，越多重置点将触发。 然后使用基于FNV算法的强哈希来产生两个重置点之间区域的哈希值。 最后将来自每个重置点FNV哈希中的字符进行串联，得到最终的模糊哈希值。

\begin{equation} \label{eq:rollinghash}
rollling hash (c_1,c_2,...,c_k) mod b =  b - 1
\end{equation}

模糊哈希算法已被应用于许多领域，因为它能够比较两个不同的文件并确定其相似性。然而，在网络流量中使用模糊哈希来识别正在传输中的相似文件时，却遇到了许多问题和挑战：

1）	无法确定文件长度：文件长度可能在文件传输结束之前无法确定。例如，在HTTP 1.0及之前版本中，content-length字段是可有可无的。在HTTP1.1及之后版本。如果是keep alive模式，则content-length和chunk字段二选一。若是非keep alive模式，则和HTTP1.0一样，content-length可有可无。文件长度虽然对于像MD5和SHA-1这样的哈希算法可有可无。但是，根据生成模糊哈希的步骤，文件长度是模糊哈希生成的关键参数之一，用于生成之前提到的触发值以生成哈希片段。

2）流式处理：为了生成适当长度的签名，模糊哈希需要完整的文件，但将网络中传输的音视频文件完全存储在硬盘上是不切实际的，这样会消耗网络中间件大量的存储资源。

3）无序输入：现在的传输技术会将文件分成片段以实现传输的高效性和灵活性，例如多线程下载，P2P文件共享和网络中继服务（例如百度云上传文件）。 如图~\ref{fig:Multi-threaded}~是典型的多线程传输场景，其中灰色块表示文件片段。在时间t3，可能出现0-3M范围内的任何文件片段。 模糊哈希算法只能从文件头或复位点进行计算，因此必须缓冲无序片段，直到收到所有前面的数据。在最坏的情况下，几乎整个文件都被缓冲，这使得内存消耗十分巨大。
\begin{figure}[!htbp]
    \centering
    %trim option's parameter order: left bottom right top
    \includegraphics[width=0.80\textwidth]{Multi-threaded}
    \bicaption{多线程传输场景}{Multi-threaded transmission scenario}
    \label{fig:Multi-threaded}
\end{figure}


4）不完全捕获：由于数据包丢失或处理错误，从网络流量捕获的文件不完整是一种常见的结果。

针对以上问题，本文采用流式模糊哈希算法以生成流量摘要。该算法采用存储中间计算结果的方法进行流式数据摘要计算，能够处理数据缺损、乱序和重叠的情况，并且采用伽罗华域（Galois Field）的乘法运算作为强哈希算法减少内存占用。

该算法有以下优点：

1）可以生成流式数据的摘要，用于网络流量的相似性检测。引入多个区间数据索引以及中间状态，以满足复杂的网络环境中数据块的无序输入。
% 流式模糊哈希处理乱序的流量输入结构图如图\ref{fig:handleout-of-order}:

% \begin{figure}[!htbp]
%     \centering
%     %trim option's parameter order: left bottom right top
%     \includegraphics[width=0.80\textwidth]{handleout-of-order}
%     \bicaption{处理乱序结构示意图}{Handling out-of-order structure}
%     \label{fig:handleout-of-order}
% \end{figure}

2）可以在内存集约的环境下工作，从而降低代理服务器、入侵检测系统等网络中间件的开销。

3）只要对流量处理一次就可以生成模糊哈希值，无需对海量的流量进行存储。

典型的流式模糊哈希值如下所示：第一个是完整数据生成的，第二个是不完整数据生成的。

\begin{lstlisting}[title=流式模糊哈希值, frame=shadowbox]
5Howh395h39h/ZF9e3Tg6S9d8CY9QXA4A5A4z48ZQdVZcuXVBAIkV2R3prfywFlL8xwY2[0:36746924]

CKStpol/OlnrBBIphB05v1E5qe[0:490021]Y8mWpVC3no6pzoEP6QCCMfZlowXRKI1JoLb[523900:950187]
\end{lstlisting}

其中方括号里的数字为数据的偏移量，例如[0:490021]表示前一段的模糊哈希值是由第0字节到第490021字节的网络流量所生成的，0是流式模糊哈希值的左偏移量，490021是流式模糊哈希值的右偏移量。

% 上文提到了，流式模糊哈希三个优点：

% 1）可以生成流式数据的摘要，用于网络流量的相似性检测。

% 2）可以在内存集约的环境下工作。 

% 3）只要对流量处理一次就可以生成模糊哈希值。

% 以上的三个优点使得流式模糊哈希可以在代理服务器中使用，以获取真实环境中的音视频流量摘要。但由于用户的观看习惯和大量的音视频网站采用流媒体技术，所以导致了缺损的不完整的音视频文件。对于完整度较低文件的相似性比较，流式模糊哈希有着一定的局限性。

% 下文将在真实环境中的代理服务器中部署流式模糊哈希算法，并计算音视频流量的摘要，以获取真实环境的数据，来帮助判断网络流中传输的音视频文件是否重复。

\section{使用基于流式模糊哈希的相似性查找系统的原因及优点}
采用流式模糊哈希值作为基准对文件的相似性进行检测时，需要使用基于流式模糊哈希的相似性查找系统。

基于流式模糊哈希的相似性查找系统采用编辑距离对流式模糊哈希值进行相似性检测，若两个流式模糊哈希值的编辑距离越小，则认为越相似。编辑距离，又称Levenshtein距离，是指两个字串之间，由一个转成另一个所需的最少编辑操作次数。许可的编辑操作包括将一个字符替换成另一个字符，插入一个字符，删除一个字符。

编辑距离计算的时间复杂度为O($n^2$)，为了减少大规模相似查找的时间，大部分的字符串查找系统采用n-gram构建索引。在\cite{Chen2008Efficient}中，作者使用n-gram作为索引，以减少候选字符串数量。n-gram是来自给定字符串的n个字符的连续序列。 为了生成n-gram，它每次从第一个字符向前移动一步字符。 例如，如果n为3，则字符串“abcde”经过处理变为“abc”，“bcd”和“cde”。 每个gram的长度可以是任何常数。

通过n-gram建立索引的一个问题是随着字符串数量的增长，一个gram的链接列表可能太长。基于流式模糊哈希的相似性查找系统通过剪枝被用来解决这个问题。在进行相似查询时，该系统将查询的模糊哈希的字符串分为几个块。为了在查询时进行加速，使用n-chunk代替n-gram。 n-chunk也是来自给定字符串的n个字符的连续序列。但是n-chunk每次从第一个角色向前移动n步。以字符串为例，如果n为3，则字符串“abcdef”经过处理变为“abc”和“def”。
% \begin{figure}[!htbp]
%     \centering
%     %trim option's parameter order: left bottom right top
%     \includegraphics[width=0.80\textwidth]{Schematicdiagram}
%     \bicaption{基于n-gram的标识匹配算法结构示意图}{Schematic diagram of n-gram-based identification matching algorithm}
%     \label{fig:Schematicdiagram}
% \end{figure}

% 然后将这些流式模糊哈希生成的gram插入哈希表中。通过n-gram建立索引的一个问题是随着字符串数量的增长，一个gram的链接列表可能太长。该查找系统通过剪枝被用来解决这个问题。

% 在进行相似查询时，系统将查询的模糊哈希的字符串分为几个块。为了在查询时进行加速，使用n-chunk代替n-gram。因为字符串的块数小于gram数。 n-chunk也是来自给定字符串的n个字符的连续序列。但是n-chunk每次从第一个角色向前移动n步。以字符串为例，如果n为3，则字符串“abcdef”将溢出为“abc”和“def”。

在此之后，将块放入哈希表中以搜索是否存在相同的块。最后评价相似度的时候使用如下公式~\ref{eq:cfds}~：

\begin{equation} \label{eq:cfds}
cfds =  10*\left(\frac{C}{chunk\_cnt}\right)\left(\frac{gram\_cnt}{left\_gram\_cnt}\right)
\end{equation}

其中C表示出现次数，chunk\_cnt表示查询的string中包含chunk的个数，gram\_cnt表示查询串的gram的个数，left\_gram\_cnt表示被比较的字符串gram个数。

该系统有以下优点：

1）专门针对流式模糊哈希值相似性查找进行设计。

2）优化了字符串匹配的查找算法，用n-chunk代替n-gram构建索引，其对800万条的查询速度可以达到2万次/秒。

\section{小结}
本章主要介绍了两个前人的工作，分别是流式模糊哈希算法和基于流式模糊哈希算法的相似性查找系统，流式模糊哈希主要是在真实环境中生成音视频流量的摘要，作为基准用以生成文件编号，这样就无需对整个网络流量中传出的音视频文件进行存储，极大的节约了存储的消耗。基于流式模糊的相似性查找系统主要用于计算两个流式模糊哈希值之间的编辑距离，评价两个文件之间的相似性。
% 其优化了字符串匹配的查找算法，其对800万条的查询速度可以达到2万次/秒。


